\scanpage{42}
\mspage{40}
% archive sheet number: 41
\deleted{a un isom. canonique}
\[
\deleted{\pi_1(S,\unsure{\xi|\zeta})\simeq\pi_1(\operatorname{Spec}k,\overline k)=\operatorname{Gal}(\overline k/k)}.
\]
\deleted{De plus, l'hom. canonique (cf.)}
\[
\deleted{\pi_1(U,\xi)\longrightarrow\pi_1(S,\xi)=\pi_0}
\]
\deleted{est surjectif, \deleted{et définit une extension dont le noyau} $I$ a la structure suivante :}

\deleted{Soit} $\unsure{E|\mathscr E}$

Quitte à diviser $E$ par son s-groupe de torsion, \inserted{(ce qui ne change pas la conclusion à prouver, ni l'hyp.)} on est ramené à examiner le cas d'un $E$ libre sur $\unsure{\mathbb Z_\ell|Z_\ell}$, l'hyp. d'action modérée s'exprimant maintenant par le fait que l'image de $I$ dans $\operatorname{Aut}(E_0)$ ($E_0=E\otimes_{\mathbb Z_\ell}\mathbb Z/\ell\mathbb Z$) est d'ordre premier à $p$. Mais on voit de suite que $I$ est la limite projective des groupes $I_\alpha$, donc l'image de $I$ dans $\operatorname{Aut}(E_0)$ est \deleted{\unsure{lim}} l'intersection de celles des $I_\alpha$, donc est égale à celle de $I_\alpha$ pour $\alpha$ grand. Donc on peut supposer \deleted{que} \inserted{$\alpha$ tel que} cette image est d'ordre premier à $p$. Donc $I'=I_\alpha$ opère \inserted{sur $E$} par l'intermédiaire d'un groupe profini \unsure{pro-fini|pro-d'ordre} de l'\deleted{entier} \inserted{ordre} transfini $\unsure{q}$ premier à $p$.

\unsure{b} On utilise maintenant le th. de structure suivant, qui est une des variantes du lemme de Abhyankar :
% slanted
\margin{devrait faire l'objet d'un n° spécial avant l'actuel n° 2}

\uline{Lemme} 8.5. Soit $A$ anneau local noeth. régulier \uline{hensélien}, $k$ son corps résiduel, $p$ l'exp. car. de $k$, \deleted{$\unsure{(x_i)_{i\in I}}$} $(x_i)_{1\leq i\leq n}$ une partie d'un syst \uline{régulier}
